% =====================================================================
%  ДОЛОО. ТӨСЛИЙН ДҮГНЭЛТ
% =====================================================================
\chapter{ТӨСЛИЙН ДҮГНЭЛТ}

\section{Төсөл хэрэгжих урьдчилсан нөхцөл, хүлээх үүрэг}

Төсөл амжилттай хэрэгжихийн тулд дараах урьдчилсан нөхцөлүүд
шийдвэрлэгдсэн байх шаардлагатай:

\begin{itemize}
  \item Хөрөнгө оруулалтын санхүүжилт (өөрийн + зээл) цагтаа
        хангагдсан байх.
  \item Төлбөрийн үйлчилгээ үзүүлэгчидтэй (QPay, SocialPay, карт
        процессинг) гэрээ байгуулсан байх.
  \item Анхдагч хэрэглэгч кофе шопуудтай хамтын ажиллагааны
        гэрээ байгуулсан байх.
  \item NFC стикер, тестийн тоног төхөөрөмж худалдаж авсан байх.
  \item Багийн бүрэн бүрэлдэхүүн ажилдаа ороод, хөгжүүлэлтийн орчин
        (cloud, repository, CI/CD \cite{githubactions}) бэлэн болсон байх.
\end{itemize}

\noindent\textbf{Хүлээх үүрэг:} Төсөл хэрэгжүүлэгч баг нь дээрх
нөхцөлүүдийг хангах, төслийн хугацаа, төсвийн дагуу гүйцэтгэх,
бүтээгдэхүүний чанарын стандартыг баримтлах үүрэг хүлээнэ. Салбар
эзэмшигч талууд нь NFC стикерийг ширээнд зөв байрлуулах, ажилтнуудаа
сургах, санал хүсэлтээ бодит цагт илгээх үүрэгтэй.

\section{Төслийн хэрэгжилтээс гарах үр дүн}

\begin{itemize}
  \item NFC стикер + PWA-д суурилсан, кассын дараалалгүй захиалга,
        төлбөрийн бүрэн системийг 6 сарын дотор бүтээн байгуулах.
  \item Эхний жилд 10 -- 20 кофе шопт бүтээгдэхүүнээ борлуулж,
        системийн тогтвортой ажиллагааг бодит орчинд батлах.
  \item Захиалгын 30 -- 40\% нь NFC урсгалаар өгөгдөх түвшинд
        хэрэглэгчдийг шилжүүлэх.
  \item Системийн таталтын ажиллагаа 99.5\%, дундаж хариу 1 секундын
        доор байх үзүүлэлтэд хүрэх.
\end{itemize}

\section{Төслийг хэрэгжүүлснээр гарах үр дүн}

\begin{itemize}
  \item \textbf{Үйлчлүүлэгчид:} Хүлээх хугацаа 40 -- 50\% багассан,
        нэг алхамаар захиалга өгөх, төлөх орчин бүрдэнэ.
  \item \textbf{Салбарууд:} Кассын ачаалал 20 -- 30\% хэмнэгдэж,
        цэс удирдах, тайлан тооцоо автоматжирана.
  \item \textbf{Төсөл хэрэгжүүлэгчид:} 3 жилийн дотор 171\%-ийн ROI,
        ойролцоогоор 35 сарын дотор хөрөнгө оруулалтаа бүрэн нөхөх
        санхүүгийн үр дүнд хүрнэ.
  \item \textbf{Салбар, эдийн засагт:} Дотоодын технологийн бүтээгдэхүүн
        нэмэгдэж, уламжлалт үйлчилгээний салбарын цифржилт хурдасна.
\end{itemize}

\section{Төслийг хэрэгжүүлснээр гарч болох эрсдлээс хамгаалах арга зам}

\begin{table}[h]
  \centering
  \caption{Эрсдэлийн үнэлгээ, хамгаалах арга зам}
  \label{tab:ch7-ersdel}
  \small
  \begin{tabular}{@{}p{0.8cm}p{3.6cm}cp{2.4cm}p{5.2cm}@{}}
    \toprule
    \textbf{№} & \textbf{Эрсдэл} & \textbf{Магадлал} & \textbf{Нөлөөлөл} &
    \textbf{Хамгаалах арга зам} \\ \midrule
    1 & Төлбөрийн үйлчилгээ үзүүлэгчийн интеграци хугацаандаа
        хоцрох & Дунд & Өндөр & Эхний үе шатнаас төлбөрийн
        документтай ажиллаж, sandbox орчинд эрт турших \\
    2 & Хэрэглэгчид NFC-ийг ойлгохгүй, QR руу буцах & Дунд & Дунд &
        Салбарын зааварчилгаа, NFC стикерийн дээр ойлгомжтой заавар
        байрлуулах, ажилтны сургалт \\
    3 & Томоохон компани ижил шийдлийг хурдан гаргах & Бага & Өндөр &
        Анхдагч хэрэглэгч кофе шопуудтай урт хугацааны хамтын
        ажиллагаа, салбарын өгөгдлийн давуу талыг бэхжүүлэх \\
    4 & Серверийн доголдол, үйлчилгээ тасалдах & Бага & Өндөр &
        Cloud-ийн өндөр хүчин чадал, автомат нөөцлөлт, мониторинг,
        offline кэш режим \\
    5 & Хувийн мэдээлэл алдагдах & Бага & Өндөр & Картын мэдээлэл
        хадгалахгүй байх, RBAC, аудитын лог, TLS, тогтмол аюулгүй
        байдлын туршилт \\
    6 & Хөгжүүлэгчийн бүрэлдэхүүн солигдох & Дунд & Дунд &
        Технологийн баримт бичгийг тогтмол хөтлөх, кодын хяналт,
        мэдлэг солилцох дотоод сургалт \\
    \bottomrule
  \end{tabular}
\end{table}

\noindent Эрсдэл бүрийг сар бүрийн төслийн хуралдаанд хянаж,
магадлал болон нөлөөллийн үнэлгээгээ шинэчилж, шаардлагатай
тохиолдолд хамгаалах арга хэмжээгээ нэмж төлөвлөнө.
